\documentclass[11pt]{article}
\usepackage[a4paper,margin=23mm]{geometry}
\usepackage{amsmath,amssymb,amsthm,mathtools,mathrsfs,slashed,bm,booktabs,microtype,xurl,hyperref}
\hypersetup{colorlinks=true,linkcolor=blue,urlcolor=blue,citecolor=blue,
pdftitle={P54 invariant attribution and nonuniform background control}}
\setlength{\emergencystretch}{2em}
\newcommand{\tr}{\operatorname{tr}}
\newcommand{\Tr}{\operatorname{Tr}}
\newcommand{\diag}{\operatorname{diag}}
\newcommand{\MS}{\overline{\mathrm{MS}}}
\newtheorem{proposition}{Proposition}
\title{Route-F P54: Exact Invariant Attribution,\\
Nonuniform Scalar Control and the Radial Fermion Kernel}
\author{P54PQ-v2 four-dimensional mainline}
\date{14 September 2026}
\begin{document}
\maketitle
\begin{abstract}
We resolve the scalar one-loop radial mass and kinetic corrections into
the actual invariants of the unchanged P54PQ-v2 action, retaining all
interference. Multihomogeneity determines the full radial Hessian jets
from one ambient Hessian per invariant, avoiding a repeated numerical
derivative scan. Fixed-propagator attribution is distinguished from
constrained parameter changes, which require new spectra, counterterms
and light-doublet tuning. This calculation concerns engineering control
of candidate backgrounds, not a fit to observed particle masses. The
radial Majorana determinant can be evaluated independently of a complete
flavor fit, but a full gauged physical stability verdict still requires
the consistently regulated momentum-dependent bosonic kernel.
\end{abstract}

\section{Unchanged theory and scope}
The action, 328 canonical real scalar coordinates and tensor
normalizations are exactly those of
\texttt{verify\_p54\_p1\_hessian\_spectrum.py}.
The reference point is
\begin{equation}
 r_0=(w,\sigma,v_s)=(1,0.1265,0.25)\omega,
 \qquad g_U=\mu_0/\omega=0.5626028899853225.
\end{equation}
Here ``original point'' means this broken background, not the symmetric
configuration $X=0$. We use the same DR/$\MS$, background-field Landau,
fixed-VEV invariant-mass counterterm prescription as P-REN1 and P-IR1.
Let $X=Rr$, with
\begin{equation}
 G=R^TR=\diag(12/5,2,1),\qquad
 H=V_0''(X),\quad T_j=\partial_{r_j}H,\quad
 U_{jk}=\partial_{r_j}\partial_{r_k}H.
\end{equation}
Thus a coordinate $\sigma\sigma$ entry must be divided by two to
obtain its value on a unit canonical $\sigma$ direction. The 38 tree
soft modes are 33 gauge Goldstones, one physical PQ mode and the four
real components of one tuned complex light doublet. There are 290
positive scalar modes. Their mass weights are recomputed for every
candidate, not copied from the original point.

\section{One Hessian per invariant: exact block Euler reconstruction}
Write the real-parameter slice of the action as
\begin{equation}
 V_0=\sum_\alpha p_\alpha I_\alpha,\qquad
 (\Phi,\Sigma,H_{10},S)=(w\Phi_0,\sigma\Sigma_0,0,v_sS_0).
\end{equation}
The real and imaginary canonical coordinates of each complex field
belong to the same scaling sector. If $I_\alpha$ has multidegree
$d_\alpha=(d_\Phi,d_\Sigma,d_H,d_S)$, the Hessian block with external
fields in sectors $A,B$ has degrees
\begin{equation}
 e_{\alpha,AB,k}=d_{\alpha,k}-\delta_{Ak}-\delta_{Bk}.
\end{equation}
Negative degrees signify a zero derivative, not a negative-power
continuation. At $H_{10}=0$, blocks with remaining $e_H>0$ also vanish.
For the other blocks and positive radial coordinates,
\begin{align}
 H^\alpha_{AB}(r')&=H^\alpha_{AB}(r_0)
     \prod_{k\in\{\Phi,\Sigma,S\}}(r'_k/r_{0k})^{e_k},\\
 (T^\alpha_j)_{AB}&=\frac{e_j}{r_j}H^\alpha_{AB},\qquad
 (U^\alpha_{jk})_{AB}=
 \frac{e_j(e_k-\delta_{jk})}{r_jr_k}H^\alpha_{AB}.
 \label{eq:Euler}
\end{align}
These identities follow by differentiating each homogeneous polynomial
block. They are exact for the fixed angular background, not a truncated
Taylor approximation or a fit. At a zero radius one must use polynomial
coefficients directly instead of dividing by that radius.

\begin{center}\small
\begin{tabular}{ll}
\toprule Couplings & Multidegree $(\Phi,\Sigma,H,S)$\\\midrule
$\mu^2;\ c;\ a,b$ & $(2,0,0,0);\ (3,0,0,0);\ (4,0,0,0)$\\
$\nu^2;\ \lambda_0,\lambda_2,\lambda_4,\lambda'_4$ & $(0,2,0,0);\ (0,4,0,0)$\\
$\alpha,\beta;\ \xi_0^2;\ \xi_1,\xi_2$ & $(2,2,0,0);\ (0,0,2,0);\ (0,0,4,0)$\\
$\xi_3;\ \gamma_1,\gamma_2,\eta_3$ & $(1,0,2,0);\ (0,2,2,0)$\\
$\eta_0,\eta_2;\ \eta_1$ & $(2,0,2,0);\ (0,3,1,0)$\\
$\mu_s^2;\ \chi_1;\ \chi_2$ & $(0,0,0,2);\ (0,0,0,4);\ (0,2,0,2)$\\
$\chi_3;\ \chi_4;\ \chi_5$ & $(2,0,0,2);\ (1,2,0,1);\ (0,0,2,2)$\\
$\chi_6;\ \chi_7$ & $(0,0,2,1);\ (1,0,2,1)$\\\bottomrule
\end{tabular}
\end{center}

For a complex coupling, $cI+c^*I^*=2c_R\Re I-2c_I\Im I$.
This checkpoint varies the real directions only. A CP-phase study
requires the independent imaginary basis and new angular-stationarity
tests. Also, a vanishing radial potential does not imply a vanishing
ambient Hessian: transverse states can couple strongly to that radius.
Conversely, $\xi_1,\xi_2$ multiply $H_{10}^4$ and have identically zero
$H,T,U$ on this entire radial slice. Changing them cannot repair this
particular one-loop radial failure, although they affect boundedness,
nonzero-$H_{10}$ backgrounds and other amplitudes.

\begin{proposition}[A structural blind spot of radial tests]
Passing every radial mass/kinetic test on $H_{10}=0$ does not imply
boundedness of the scalar action.
\end{proposition}
\begin{proof}
Set $\Phi=\Sigma=S=0$. The quartic potential is
\begin{equation}
 V_4=\xi_1(H^\dagger H)^2+\xi_2|H^TH|^2.
\end{equation}
For a complex null vector $H=t(1,i,0,\ldots)/\sqrt2$, one has
$H^TH=0$ and $V_4=\xi_1t^4$. Thus $\xi_1<0$ makes the potential
unbounded below without changing any of the radial $H,T,U$ computed
above. A real $H$ similarly requires $\xi_1+\xi_2\ge0$. Since
$0\le |H^TH|^2/(H^\dagger H)^2\le1$, these two conditions characterize
nonnegativity of this pure-$H$ quartic. They are only necessary for the
full action; equality still requires lower-degree terms on flat rays,
and mixed-field directions have not been bounded.
\end{proof}
Both necessary inequalities are now explicit rejection gates. The
global mixed-quartic and competing-vacuum problems remain separate;
additional radial decimal accuracy cannot replace either one.

\section{Attribution with all interference retained}
Diagonalize only the complete action:
\begin{equation}
 Q^THQ=\diag(m_i^2),\qquad
 \widetilde T^\alpha_j=Q^TT^\alpha_jQ,\qquad
 \widetilde U^\alpha_{jk}=Q^TU^\alpha_{jk}Q.
\end{equation}
Here component jets include the actual coefficient $p_\alpha$.
Individual invariants need not have positive Hessians and do not define
separate physical determinants. Fix $Q,m_i^2$ and the hard projector
only for the following attribution. Define
\begin{align}
 A_i&=m_i^2[\ln(m_i^2/\mu_0^2)-1]\quad (i\text{ hard}),\quad A_i=0\ (i\text{ soft}),\\
 L(a,b;s)&=\int_0^1\ln\frac{(1-x)a+xb+x(1-x)s}{\mu_0^2}\,dx,\quad s=p_E^2,\\
 W(a,b)&=\int_0^1\frac{x(1-x)}{(1-x)a+xb}\,dx.
\end{align}
In the local hard kernels below, at least one internal line is hard.
Hard--soft terms are retained. Pure soft--soft terms instead contain
$L(0,0;s)=\ln(s/\mu_0^2)-2$ and are not a local kinetic coefficient.

The scalar tadpole and common finite counterterm for component $\alpha$
are linear:
\begin{align}
 t^\alpha_j&=\frac{1}{32\pi^2}\sum_i A_i(\widetilde T^\alpha_j)_{ii},\\
 \mathcal A(r)&=-\diag(12w/5,\sigma,v_s),\qquad
 \Delta p^\alpha=-\mathcal A(r)^{-1}t^\alpha,\\
 R^TH_{\mathrm{CT}}^\alpha R
 &=-\diag(t^\alpha_w/w,t^\alpha_\sigma/\sigma,t^\alpha_v/v_s).
\end{align}
Hence define
\begin{equation}
 C^\alpha_{jk}=\frac{1}{32\pi^2}\sum_i A_i
 (\widetilde U^\alpha_{jk})_{ii}
 +(R^TH_{\mathrm{CT}}^\alpha R)_{jk}.
\end{equation}
The ordered bubble and kinetic interference blocks are
\begin{align}
 B^{\alpha\beta}_{jk}&=\frac{1}{32\pi^2}\sum_{il}'
 (\widetilde T^\alpha_j)_{il}(\widetilde T^\beta_k)_{li}
 L(m_i^2,m_l^2;0),\\
 Z^{\alpha\beta}_{jk}&=\frac{1}{32\pi^2}\sum_{il}'
 (\widetilde T^\alpha_j)_{il}(\widetilde T^\beta_k)_{li}
 W(m_i^2,m_l^2).
\end{align}
The prime excludes only two-soft-line pairs. All ordered blocks are
saved, including $\alpha\ne\beta$. Their exact sums give
\begin{equation}
 H_{S+\mathrm{CT}}=\sum_\alpha C^\alpha+\sum_{\alpha\beta}B^{\alpha\beta},
 \qquad \Delta Z_S=\sum_{\alpha\beta}Z^{\alpha\beta}.
\end{equation}
The joint $(\alpha,j)$ kinetic Gram matrix is positive semidefinite,
but a cross block need not be positive. The signed equal-share ledger
\begin{align}
 \mathscr M^\alpha&=C^\alpha+\tfrac12\sum_\beta
 (B^{\alpha\beta}+B^{\beta\alpha}),\\
 \mathscr Z^\alpha&=\tfrac12\sum_\beta
 (Z^{\alpha\beta}+Z^{\beta\alpha})
\end{align}
is a specified accounting convention, not a separately measurable
``contribution percentage.'' Its sums reconstruct the full matrices.

\subsection{Why the ledger is not a parameter derivative}
If $\kappa_\alpha$ scales only vertices at fixed propagators,
\begin{equation}
 F(\kappa)=\sum_\alpha\kappa_\alpha C^\alpha+
 \sum_{\alpha\beta}\kappa_\alpha\kappa_\beta B^{\alpha\beta},\qquad
 \partial_{\kappa_\alpha}F\big|_1
 =C^\alpha+\sum_\beta(B^{\alpha\beta}+B^{\beta\alpha}).
\end{equation}
Even this frozen derivative differs from the signed allocation. Actual
parameter changes also differentiate masses, eigenprojectors, loop
weights and the stationarity/light-doublet constraints. The three
dependent mass terms illustrate the distinction: their cubic/quartic
jets vanish, but they still change all propagators.

\section{Actual invariant ledger at the original point}
The exact reconstruction, random independently reweighted vertex tests,
full joint kinetic Gram positivity and source comparisons pass 445/445
checks. The first run required 27 unit-coefficient Hessians (the two
$H^4$ zeros are algebraic) and one new independent full-action radial
point. Subsequent runs use those caches without new evaluations.
For the canonical $\sigma$ direction $e_\sigma/\sqrt2$ and the
$G$-normalized leading kinetic eigenvector $z$, the signed ledger is
\begin{center}
\begin{tabular}{lrrr}
\toprule Invariant & $\mathscr M_{\sigma\sigma}/(2\omega^2)$ &
$z^T\mathscr Z z$ & $z^TZ^{\alpha\alpha}z$\\\midrule
$\lambda'_4$ & $-0.29460220$ & $0.89223875$ & $0.37939853$\\
$\lambda_2$ & $-0.25234309$ & $0.85449316$ & $0.29839039$\\
$\lambda_4$ & $-0.19391061$ & $0.70984670$ & $0.21519366$\\
$\lambda_0$ & $-0.02977324$ & $0.09307175$ & $0.00598139$\\
$\alpha$ & $-0.00000790$ & $0.06652577$ & $0.00226620$\\
$\beta$ & $0.00000601$ & $-0.01376498$ & $0.00206069$\\\bottomrule
\end{tabular}
\end{center}
The three transverse $\Sigma$ quartics account for $2.45657860$ of
the leading $2.61031989$ kinetic insertion (94.11\% in this specified
allocation), and $-0.74085590\,\omega^2$ of the scalar-plus-CT canonical
$\sigma$ curvature (96.13\%). These percentages are not invariant
observables. In particular all kinetic self terms sum to only
$0.90706619$; discarding interference misses $1.70325369$.

The radial restriction of the actual tensor action contains
$\lambda_0\sigma^4/4$, but not these three transverse quartics.
At a stationary background
\begin{equation}
 (H_0)_{\sigma\sigma}=2\lambda_0\sigma^2,
 \qquad R^TH^{\lambda_2,\lambda_4,\lambda'_4}R=0.
\end{equation}
The latter projected numerical norms are below $2\times10^{-16}$ at
the original card. Thus transverse quartics can be weakened without
simultaneously weakening this tree radial restoring term. This is a
directional reason to test nonuniform changes, not a guarantee of
success: the propagator spectrum still changes.

\section{Constrained nonuniform re-evaluation}
At each separately declared candidate card, assemble the full polynomial
$H,T,U$ by Eq.~\eqref{eq:Euler}. Re-solve the three mass parameters by
\begin{equation}
 \delta p=-\mathcal A(r)^{-1}\nabla_r V_0,
 \qquad H\mapsto H+H_{\mathrm{mass}}(\delta p),
\end{equation}
and tune $\xi_0^2$ to the first physical eigenvalue crossing, after
removing the 34 symmetry directions. Verify the four extra zero modes
have $(C_2^{SU(3)},C_2^{SU(2)},Y^2)=(0,3/4,1/4)$ and all other scalar
tree modes are positive. Then recompute every scalar/vector mass,
radial jet contraction and the common finite tadpole shift.

The engineering tests are
\begin{equation}
 H_b^{\rm hard}>0,\qquad
 \rho_H=\max|\lambda(H_b^{\rm hard}-H_0,H_0)|<0.3,
 \qquad \rho_Z=\lambda_{\max}(\Delta Z_S,G)<0.3.
\end{equation}
The margin $0.3$ is declared, not a theorem bounding omitted loops.
Passing these radial tests does not prove full nonradial stability,
boundedness from below, a physical pole, or a valid unification fit.
No observed mass or flavor data are used, and no default matching input
is overwritten.

\subsection{First-stage cards: weakening quartics alone is insufficient}
The six declared cards use
$(\lambda_2,\lambda_4,\lambda'_4)\mapsto
\tau(\lambda_2,\lambda_4,\lambda'_4)$.
Except as labelled, $r=r_0$ and all other independent coefficients are
unchanged. L2 doubles $\lambda_0$; X multiplies $\chi_4$ by $0.03$;
S30 also changes $\sigma/\omega$ to $0.3$.
\begin{center}\small
\begin{tabular}{lrrrr}
\toprule Card & $\tau$ & $\min\lambda(H_b^{\rm hard},G)/\omega^2$
& $\rho_H$ & $\rho_Z$\\\midrule
T30 & $0.30$ & $-0.434370$ & $32.3860$ & $0.902652$\\
T10 & $0.10$ & $-0.106794$ & $8.59012$ & $0.713427$\\
T03 & $0.03$ & $-0.026376$ & $2.84877$ & $1.099954$\\
T03 L2 & $0.03$ & $-0.067918$ & $3.31182$ & $1.370839$\\
T03 X & $0.03$ & $-0.026349$ & $2.84547$ & $1.099232$\\
T03 X S30 & $0.03$ & $-0.052148$ & $1.66116$ & $0.632098$\\\bottomrule
\end{tabular}
\end{center}
All six fail. Their actual action reconstruction, necessary BFB negative
controls and 80 independent bubble integrals pass 183/183 numerical
checks; rejection is a physical-control result, not a failed code test.
At T03 X S30, a separate evaluation of the final retuned full action,
three coordinate-direction jets and a mixed-direction jet confirms
that the invariant assembly is not merely an assumed input identity.

The nonmonotonic kinetic insertion has a concrete spectral origin.
Between T10 and T03 the positive scalar gap drops from
$0.00883111$ to $0.00267622\,\omega^2$, while the leading kinetic
eigenvector's canonical $\Phi$ fraction rises from $0.667676$ to
$0.824065$. In this new mass basis the signed $\alpha$ allocation
rises from $0.437312$ to $1.081219$ (the accompanying $\beta$
interference is negative). Thus the original leading-$\Sigma$
quartic problem is replaced by a large mixed-coupling response in the
$\Phi$ direction. Merely shrinking a vertex does not uniformly shrink
a loop with a simultaneously closing propagator gap.

For example, any positive hard squared-mass gap $\Delta$ gives
\begin{equation}
 W(a,b)\le\frac{1}{6\Delta}\quad(a,b\ge\Delta),\qquad
 W(0,b)\le\frac{1}{2\Delta}\quad(b\ge\Delta).
\end{equation}
Therefore a useful control estimate must constrain both the vertex
norm and the inverse gap. A fixed-parameter derivative of a mass is
not a derivative along the reanchored stationary-mass curve; confusing
the two would incorrectly remove precisely the response measured here.

\subsection{Second-stage joint mixed-coupling probes}
After identifying this response, two additional coarse cards were
declared, not optimized: $\lambda_0=0.2$ as an absolute value;
$(\lambda_2,\lambda_4,\lambda'_4,\chi_4)$ are multiplied by $0.01$;
$(\alpha,\beta,\chi_2)$ by $0.1$. All other independent coefficients
are retained, with $w=1$, $v_s=0.25$ and $\sigma=0.3$ or $0.5$.
Dependent masses and the tree doublet are re-solved as above.
\begin{center}
\begin{tabular}{rrrr}
\toprule $\sigma/\omega$ & $\min\lambda(H_b^{\rm hard},G)/\omega^2$
& $\rho_H$ & $\rho_Z$\\\midrule
$0.3$ & $-0.00317747$ & $1.176951$ & $1.131599$\\
$0.5$ & $0.02299103$ & $0.538986$ & $1.899032$\\\bottomrule
\end{tabular}
\end{center}
The 36 checks pass with zero new Hessian evaluations, but neither card
passes the declared control gates. The second is an explicit example
of positive hard curvature without small loop insertions. Its leading
kinetic vector is 99.9798\% in the canonical $\Phi$ direction. Neither
is retained, so no new independent full-action Hessian calculation is
spent on these two rejected cards; their assembly uses the checked
invariant library. This is the end of this bounded eight-card test,
not the start of an unrestricted optimization.

\subsection{A conditional inverse-gap obstruction, not a decimal fit}
On the ray changing only the three transverse quartics at fixed $r_0$,
two mutually orthogonal scalar eigenspaces have numerically verified
projector identities (the following numerical matrices use $\omega=1$)
\begin{align}
 P_{50}H P_{50}&=12\tau\sigma^2P_{50},&
 P_{36}H P_{36}&=\frac{58}{5}\tau\sigma^2P_{36},\\
 P_j T_w P_j&=\frac9{25}P_j,& j&=50,36.
\end{align}
The projectors are pure $\Sigma$, with ranks indicated by their
subscripts. Their entries admit exact integer-over-eight projector
certificates. The action-to-projector identities above are checked
against the actual tensor action, but a symbolic or interval
certificate for those action identities is not yet supplied. The
following implication is therefore an exact conditional bound, not
an unconditional theorem about the entire Spin(10) action.

The separately tested invariant coefficients give the useful mass
branch formula
\begin{equation}
 m_\kappa^2=-\frac{\nu^2}{2}
 +(\lambda_0/2+\tau\kappa)\sigma^2
 +\frac65(\alpha-\beta)w^2+60\chi_2v_s^2,
 \quad (N,\kappa)=(50,12),(36,58/5).
\end{equation}
For general transverse quartics the two $\kappa$ combinations are
$4\lambda_2+4\lambda_4+16\lambda'_4$ and
$6\lambda_2+6\lambda_4+8\lambda'_4$ before multiplying by $\tau$.
The stationary $\nu^2$ cancels all but $\tau\kappa\sigma^2$ in the
mass, but \emph{does not} cancel the fixed-Lagrangian mass derivative
\begin{equation}
 b_\kappa=\nabla_r m_\kappa^2
 =\left(\frac{12}{5}(\alpha-\beta)w,
 (\lambda_0+2\tau\kappa)\sigma,120\chi_2v_s\right).
\end{equation}
This is the concrete reason that a light stationary spectrum can have
large couplings to background fluctuations.

\begin{proposition}[Conditional scalar kinetic lower bound]
Assuming these reducing-projector identities and a well-defined
positive hard complement, the scalar hard kinetic matrix obeys
\begin{align}
 (\Delta Z_S)_{ww}&\ge
 \frac{(9/25)^2}{192\pi^2\tau\sigma^2}
 \left(\frac{50}{12}+\frac{36}{58/5}\right)
 =\frac{0.03107164552}{\tau},\\
 \rho_Z&\ge\frac{(\Delta Z_S)_{ww}}{G_{ww}}
 \ge\frac{0.01294651897}{\tau}.
\end{align}
\end{proposition}
\begin{proof}
Each of the rank-$N$ degenerate blocks contributes
$N(9/25)^2W(m^2,m^2)/(32\pi^2)$ to this diagonal kinetic entry.
Since $W(m^2,m^2)=1/(6m^2)$ and every omitted scalar-pair square has
nonnegative weight, retaining the two blocks is a lower bound.
The generalized Rayleigh quotient in the $w$ direction supplies
the second inequality.
\end{proof}
Consequently $\tau<0.0431551$ cannot satisfy the declared $\rho_Z<0.3$
scalar gate on this ray under the stated identities. If another mode
becomes tachyonic first, the candidate already fails its tree gate.
The coefficient is obtained from the block masses and ranks, not a
fit to the sampled $\rho_Z$ values. This says nothing by itself about
the complete gauge-dependent kinetic matrix or physical pole stability.
It does show why continued weakening of these quartics alone is not
a reliable repair strategy.

One can strengthen the bound without any further spectrum scan. Put
$b_0=(0.36,0.11385,0.03)$ and
$z_0=G^{-1}b_0/\sqrt{b_0^TG^{-1}b_0}$. Since all these components and
$\tau$ are positive, $z_0^Tb_\kappa\ge\sqrt{b_0^TG^{-1}b_0}$.
Applying this single fixed witness to the sum of the two actual
rank-one Gram blocks yields
\begin{equation}
 \rho_Z\ge\frac{b_0^TG^{-1}b_0}{192\pi^2\tau\sigma^2}
 \left(\frac{50}{12}+\frac{36}{58/5}\right)
 =\frac{0.01471609503}{\tau}.
\end{equation}
Thus $\tau<0.0490537$ already fails the $0.3$ scalar margin under the
same verified-but-not-symbolically-certified action identities.
This witness proof does not assume the generally false matrix inequality
$b_\kappa b_\kappa^T\succeq b_0b_0^T$.

\section{Radial fermion kernel without a global family fit}
The already checked Yukawa tensor map on this same radial family is
$M_F=\sigma F_M$, with no $w$ or $v_s$ vertex. A Takagi transformation
therefore simultaneously diagonalizes the radial mass and its vertex:
$M_i=\sigma f_i$, $y_i=\partial_\sigma M_i=f_i$, for three Majorana
masses. The remaining family spinor directions have zero radial mass
and vertex. The three $f_i$ are inputs, not fitted numbers.

For one four-component Majorana fermion, expansion of
$-\tfrac12\Tr\log(\slashed\partial+M+y\delta\sigma)$ gives
\begin{align}
 \Pi_F^{\rm raw}(s)&=2y^2\int_k
 \frac{M^2-k\cdot(k+p)}{(k^2+M^2)((k+p)^2+M^2)}\\
 &=-\frac{y^2}{16\pi^2}\left[
 (4M^2+s)L(M^2,M^2;s)+2M^2(\ln(M^2/\mu_0^2)-1)\right].
\end{align}
The factor one half is Majorana counting; a Dirac field with the same
couplings would contribute twice as much. The same determinant gives
\begin{equation}
 V_F=-\frac{M^4}{32\pi^2}[\ln(M^2/\mu_0^2)-3/2],\qquad
 t_\sigma^F/\sigma=-\frac{2y^2M^2}{16\pi^2}[\ln(M^2/\mu_0^2)-1].
\end{equation}
Subtracting $t_\sigma^F/\sigma$ by the same invariant finite mass
counterterm, with no momentum-dependent tadpole subtraction, yields
\begin{equation}
 \boxed{\Pi^{F,\mathrm{FV}}_{\sigma\sigma}(s)=
 -\frac{1}{16\pi^2}\sum_{i=1}^3\frac{M_i^2}{\sigma^2}
 (4M_i^2+s)L(M_i^2,M_i^2;s).}
\end{equation}
All other radial entries are zero. Its local finite kinetic term is
\begin{equation}
 \boxed{\Delta Z^F_{\sigma\sigma}=-\frac{1}{16\pi^2}
 \sum_{i=1}^3\frac{M_i^2}{\sigma^2}
 [\ln(M_i^2/\mu_0^2)+2/3].}
\end{equation}
The canonical $\sigma$ coefficient is half this expression because
$G_{\sigma\sigma}=2$. Both expressions have the continuous zero-Yukawa
limit; no artificial $\log 0$ is evaluated. At zero momentum,
\begin{equation}
 \Pi^{F,\mathrm{FV}}_{\sigma\sigma}(0)
 \le\frac{6\mu_0^4}{e\,16\pi^2\sigma^2}
 =0.0875112373\,\omega^2,
\end{equation}
recovering P-REN1 by maximizing $-x^2\ln(x/\mu_0^2)$ at
$x=\mu_0^2e^{-1/2}$. This bound is not asserted at nonzero momentum.

The independent verifier compares the closed form with the shifted
spin-trace numerator integral
\begin{equation}
 \Pi_F^{\rm raw}(s)=\frac{2y^2}{16\pi^2}\int_0^1
 D(x)[1-3\ln(D(x)/\mu_0^2)]\,dx,\quad D(x)=M^2+x(1-x)s,
\end{equation}
followed by its tadpole subtraction. Five synthetic mass cards at five
momenta, independent kinetic finite differences, zero-mass limits,
the old curvature function and its exact cap pass 56/56 tests together
with the gauge tensor checks below. These are regression inputs,
not observed or fitted families. A full nonradial fermion kernel still
requires the complex $h_{\rm raw},f_{\rm raw}$ matrices.

\section{What the gauge spectrum misses at nonzero momentum}
To make the background-field label operational, choose the Euclidean
quadratic gauge fixing, with real antisymmetric generators $T^a$,
\begin{equation}
 D_\mu X=\partial_\mu X+g a_\mu^aT^aX,\quad
 X=\bar X+\eta,\quad v_a=T^a\bar X,\quad
 F^a=\partial_\mu a_\mu^a+\xi g v_a^T\eta.
\end{equation}
The quadratic bosonic operator and ghost operator are
\begin{align}
 Q_B&=\begin{pmatrix}Q_{SS}&W^T\\W&Q_{VV}\end{pmatrix},\qquad
 Q_{SS}=-\partial^2+V''+\xi g^2\sum_a v_av_a^T,\\
 (Q_{VV})^{ab}_{\mu\nu}&=
 [-\partial^2\delta_{\mu\nu}+(1-\xi^{-1})\partial_\mu\partial_\nu]\delta^{ab}
 +M_{ab}^2\delta_{\mu\nu},\\
 M_{ab}^2&=g^2v_a^Tv_b,\qquad
 W^a_{\mu I}=-2g(T^a\partial_\mu\bar X)_I,\qquad
 Q_c=-\partial^2+\xi M^2.
\end{align}
The mixed vertex follows directly from integrating the gauge-fixing
cross term by parts. The scalar kinetic cross term is
$g a_\mu[v_a^T\partial_\mu\eta-(\partial_\mu v_a)^T\eta]$;
gauge fixing cancels the first term and doubles the second. Consequently
$W=0$ for a constant background, but not for a radial external momentum.
With $d_{ar}=T^aR_r$, the actual additional tensors are
\begin{align}
 B_{r,ab}&=g^2(d_{ar}^Tv_b+v_a^Td_{br}),\\
 C_{rs,ab}&=g^2(d_{ar}^Td_{bs}+d_{as}^Td_{br}),\qquad
 W^a_{r,\mu I}(p)=-2igp_\mu(d_{ar})_I.
\end{align}
Their numerical values and the identities
$r_rB_r=2M^2$, $r_rr_sC_{rs}=2M^2$ are saved. The generalized eigenvalues
of $g^2\sum_a d_{ar}^Td_{as}$ are $(0,3.1652201,3.9565251)$.
This is a vertex Gram matrix, \emph{not} a loop kinetic correction.

The exact diagram generator is
\begin{equation}
 \Gamma_1^B=\tfrac12\Tr\log Q_B-\Tr\log Q_c+\Gamma_{\rm CT},\quad
 \Pi^B_{rs}=\tfrac12\Tr[Q_B^{-1}Q_{B,rs}
 -Q_B^{-1}Q_{B,r}Q_B^{-1}Q_{B,s}]+\Pi^c_{rs}+\Pi^{\rm CT}_{rs}.
\end{equation}
The remaining integrals include the two-vector bubble, vector seagull
and scalar--vector derivative bubble. In strict Landau,
$D_{a,\mu\nu}(k)=(\delta_{\mu\nu}-k_\mu k_\nu/k^2)/(k^2+m_a^2)$,
and the latter contains
\begin{equation}
 \Pi^{SV}_{rs}(p)=-4g^2\sum_{a,k}(d_{ar}\cdot u_k)(d_{as}\cdot u_k)
 \int_q\frac{p_\mu D_{a,\mu\nu}(q)p_\nu}{(q+p)^2+m_k^2}.
\end{equation}
Gauge generators must be rotated by the vector mass eigenbasis here.
This dimensionally regulated integral has not yet been evaluated in the
complete P54 prescription. Replacing the Lorentz trace by the integer
three before pole subtraction loses finite $\epsilon\times1/\epsilon$
terms. Similarly, direct background-scalar ghost vertices are proportional
to $\xi$ in this explicit fixing and can vanish in the strict Landau
limit, but the finite-$\xi$/Nielsen check, and ghost terms in general
matching, remain necessary. No guessed ghost counterterm is introduced.

\section{Physical interpretation and outstanding work}
Goldstone resummation of tadpoles alone does not make the second
derivative of the Landau effective potential finite. External momentum
is required; beyond one loop it must be combined with the appropriate
resummation \cite{BG}. The cited two-loop mass construction specializes
to the gaugeless limit and cannot be substituted wholesale for the
present gauged theory. The VEV/tadpole convention must also be kept
consistent with gauge dependence \cite{DL}.

A physical stability claim requires the complete BRST-consistent
inverse propagator and physical projection, including scalar--vector
mixing when allowed, vector/longitudinal/ghost contributions, fermions,
counterterms, analytic continuation and the relevant Ward/Nielsen
identities. Neither a negative partial hard curvature nor an isolated
root of a scalar-only kernel is such a claim. Full Wilson/box matching,
lower gauge/ghost/Yukawa finite matching and the complete finite
seesaw including $C_{HN}$ remain open. A flavor fit is not promoted.

\begin{thebibliography}{9}
\bibitem{BG} J. Braathen and M. D. Goodsell,
``Avoiding the Goldstone Boson Catastrophe in general renormalisable
field theories at two loops,'' \href{https://arxiv.org/abs/1609.06977}{arXiv:1609.06977}.
\bibitem{DL} V. D\={u}d\.{e}nas and M. L\"oschner,
``Vacuum expectation value renormalization in the Standard Model and
beyond,'' \href{https://arxiv.org/abs/2010.15076}{arXiv:2010.15076}.
\end{thebibliography}
\end{document}
